\documentclass[11pt]{article}
\usepackage[margin=1in]{geometry}
\usepackage{amsmath,amssymb}
\usepackage{booktabs}
\usepackage{calc}
\usepackage{longtable}
\usepackage{array}
\usepackage{graphicx}
\usepackage{float}
\usepackage{caption}
\usepackage{setspace}
\usepackage{parskip}
\usepackage{hyperref}
\usepackage{xcolor}
\usepackage{enumitem}
\usepackage{titlesec}

\setstretch{1.08}
\setlength{\parindent}{0pt}
\setlength{\parskip}{0.6em}
\setlist[itemize]{topsep=3pt,itemsep=2pt,leftmargin=1.5em}
\setlist[enumerate]{topsep=3pt,itemsep=2pt,leftmargin=1.5em}

\hypersetup{
  colorlinks=true,
  linkcolor=black,
  urlcolor=blue,
  citecolor=black
}

\titleformat{\section}{\large\bfseries}{\thesection.}{0.5em}{}
\titleformat{\subsection}{\normalsize\bfseries}{\thesubsection}{0.5em}{}

\begin{document}

\begin{center}
{\LARGE Final Report: Reinforcement Learning for the Codenames Spymaster
Task \par}
\vspace{0.75em}
\end{center}

\vspace{1em}

\subsection{Introduction}\label{introduction}

Our project studies the spymaster role in \emph{Codenames} as a
reinforcement learning problem. In the game, the spymaster sees the
hidden roles of all words on the board and must give a one-word clue and
a number that help a teammate guess friendly words while avoiding
opponent words, neutral words, and especially the assassin. This is a
good RL setting because each action requires balancing semantic
similarity, long-horizon planning, and risk. A clue that is too broad
may accidentally point to a dangerous word, while a clue that is too
narrow may fail to make enough progress toward winning.

The main goal of the project was to build a learning-based spymaster
that uses the semantic structure of the board rather than relying only
on hand-designed rules. Following the original proposal, we focused on
four ideas: reward shaping, behavioral cloning, Soft Actor-Critic (SAC),
and Hindsight Experience Replay (HER). Our final project includes all
four. We built a goal-conditioned version of the task, a strong greedy
semantic baseline, demonstration generation for behavioral cloning, and
a SAC training pipeline with optional HER and ablations. Because our
final large-scale experiments are still being completed, the final
result tables in this draft are marked \texttt{TBD}.

More broadly, this project sits at an interesting boundary between
reinforcement learning and language-like reasoning. Codenames is not a
standard control task in which the agent chooses from a small set of
physical actions. Instead, the spymaster must choose a clue that is
semantically related to some words on the board and semantically
separated from others. That makes the problem a useful testbed for
whether RL methods from class can be adapted to settings where
representation quality matters as much as exploration and optimization.

Our original hypothesis was that a purely model-free RL approach would
struggle unless we injected more structure into the problem. In
particular, we expected that three additions would matter: a semantic
representation of clues and board words, a shaped reward that gives
useful turn-level feedback, and a demonstration-based warm start so the
agent does not begin from completely random behavior. The final project
was designed to test that hypothesis.

\subsection{Problem Statement}\label{problem-statement}

The task is to learn a policy for the spymaster in \emph{Codenames}. At
each turn, the spymaster sees the full board with hidden roles and must
output a one-word clue and a positive integer count. A teammate then
uses that clue to guess board words. The goal is to reveal all friendly
words before hitting the assassin, while avoiding opponent and neutral
words whenever possible.

From an RL perspective, this creates several challenges. First, the
action space is unusual because clues are words, not ordinary
low-dimensional controls. Second, the reward is sparse at the game
level: the agent may make several reasonable clueing decisions before
eventually losing. Third, clue quality depends on semantic
relationships, so the representation of words matters directly to the
learning problem. Finally, there is a tradeoff between ambition and
safety. A clue that aims at many friendly words can shorten the game,
but it also increases the chance of pointing the guesser toward a bad
word.

We therefore formulated the problem as a semantic decision task. Instead
of trying to model every aspect of open-ended language generation, we
treated clueing as selecting a point in a semantic space and then
decoding that point into a legal clue. This preserved the spirit of the
game while making the problem compatible with the RL tools we wanted to
use.

\subsection{What We Did}\label{what-we-did}

We modeled the Codenames spymaster as an agent that observes the full
hidden board and outputs a clue. The environment then simulates a fixed
teammate guesser, evaluates the outcome of that clue, and returns reward
to the spymaster.

Each episode uses a 5x5 board drawn from a fixed 400-word Codenames
vocabulary in \texttt{data/raw/codenames\_words.txt}. In the default
setup, the board contains 8 friendly words, 8 opponent words, 8 neutral
words, and 1 assassin. The observation includes semantic embeddings for
the board words, a pairwise board similarity matrix, role labels, a mask
for unrevealed words, and a goal representation. The environment is
goal-conditioned: on each turn it chooses a small subset of remaining
friendly words as the current target, and the agent is rewarded for
giving a clue that helps reveal that subset safely.

The agent's action is represented as a continuous vector. Most of the
vector corresponds to a clue embedding, and the final scalar corresponds
to the clue count. The environment decodes that vector into the nearest
legal clue word from a filtered clue vocabulary and converts the count
into an integer between 1 and 9. This lets us use SAC, which is
naturally designed for continuous actions, while still producing valid
game actions.

The final implementation uses sentence-transformer embeddings from
\texttt{all-MiniLM-L6-v2}. The clue vocabulary is drawn from a curated
list, filtered to single-word clues, and screened to remove obvious
illegal overlaps with board words. In the base configuration, the clue
vocabulary is capped at 8000 candidates.

This design means that the learned agent is not generating language in
the full open-ended sense. Instead, it is solving a structured semantic
control problem. For our purposes, that was a feature rather than a
limitation: it let us focus on the central RL question of whether the
agent could learn to give better clues over time given the board
structure, reward function, and demonstrations.

\subsection{Approach}\label{approach}

Our approach combines a strong heuristic baseline with off-policy
reinforcement learning. The baseline searches over legal clue words and
chooses the clue and count that maximize a cosine-margin objective: the
clue should be close to the intended friendly targets and far from the
most dangerous alternative. This baseline serves two purposes. First, it
is the main comparison point for the learning-based agent. Second, it
generates demonstration trajectories for behavioral cloning.

The learning agent uses Stable-Baselines3 SAC with a multi-input policy.
Before RL fine-tuning, we optionally pretrain the actor on
demonstrations from the greedy baseline using behavioral cloning. The BC
stage minimizes mean-squared error between the actor's output action
vector and the demonstration action vector. We can also seed the replay
buffer with those demonstrations so that early exploration starts from a
more meaningful set of transitions.

HER is used because the task has sparse elements even after shaping. A
clue may be partially useful even if it does not fully solve the turn,
and a game may contain many partially successful states before a
complete win. In our environment, each turn has a desired goal
consisting of a coherent subset of remaining friendly words. HER allows
transitions to be relabeled in terms of achieved goals, which is a
natural fit for this kind of partial progress.

Reward shaping was central to our method. The reward combines three
signals. First, we use a semantic clue-margin term that measures how
well the clue separates intended friendly words from dangerous words.
Second, we add penalties when the guesser reveals an opponent, neutral,
or assassin word. Third, we add a turn penalty when the desired goal
subset is not achieved. The goal was to avoid a purely sparse ``win or
lose'' signal, which would make the problem much harder to learn from.

Conceptually, the final method can be summarized as follows. We start
with a semantic representation of the board and clue space. We define a
strong non-learning baseline in that same semantic space. We use the
baseline both as a comparator and as a source of demonstrations. We then
train a continuous-action RL policy with SAC, optionally warm-started by
behavioral cloning and assisted by HER and reward shaping. This made the
project much more than a direct application of one algorithm; it became
an exercise in structuring the problem so that the algorithm had a
realistic chance to learn.

\subsection{Method Summary Table}\label{method-summary-table}

\begin{longtable}[]{@{}
  >{\raggedright\arraybackslash}p{(\columnwidth - 4\tabcolsep) * \real{0.3333}}
  >{\raggedright\arraybackslash}p{(\columnwidth - 4\tabcolsep) * \real{0.3333}}
  >{\raggedright\arraybackslash}p{(\columnwidth - 4\tabcolsep) * \real{0.3333}}@{}}
\toprule\noalign{}
\begin{minipage}[b]{\linewidth}\raggedright
Component
\end{minipage} & \begin{minipage}[b]{\linewidth}\raggedright
Final choice
\end{minipage} & \begin{minipage}[b]{\linewidth}\raggedright
Why we used it
\end{minipage} \\
\midrule\noalign{}
\endhead
\bottomrule\noalign{}
\endlastfoot
Task formulation & Goal-conditioned spymaster-only RL & Simplifies
evaluation and isolates clue quality \\
Board setup & 5x5 Codenames board & Matches the standard game format
closely enough for the project \\
Representation & Sentence-transformer embeddings & Gives a usable
semantic space for both clues and board words \\
Action space & Continuous clue embedding + scalar count & Makes SAC
applicable while preserving semantic flexibility \\
Guesser & Fixed greedy teammate & Stabilizes evaluation and reduces
multi-agent complexity \\
Baseline & Greedy cosine-margin clue generator & Strong semantic
baseline and source of demonstrations \\
RL algorithm & SAC & Off-policy continuous-control method covered in
class and well suited to the action design \\
Replay strategy & HER & Helps with partial progress and sparse-goal
structure \\
Pretraining & Behavioral cloning from greedy demonstrations & Improves
initialization over random exploration \\
Reward & Margin shaping + bad-guess penalties + goal penalty & Gives
denser signal than pure win/loss reward \\
\end{longtable}

\subsection{Related Ideas That Influenced the
Project}\label{related-ideas-that-influenced-the-project}

There was not space in our presentation for a full related-work review,
but a few ideas heavily influenced the final design. SAC and HER
directly motivated the RL side of the project, since we wanted an
off-policy algorithm that works well with continuous actions and a
replay strategy that can help on sparse-goal problems. Sentence-BERT
style embeddings influenced the semantic representation, since the core
difficulty in Codenames is measuring whether a clue groups the right
words and separates them from the wrong ones. More broadly, the project
was influenced by prior work on embedding-based semantic reasoning in
word games and language-adjacent control tasks, even though our final
system is purpose-built for this course project rather than a direct
reproduction of a single paper.

The most important influence was not a single prior result, but a
pattern from recent ML work: when the task has language-like structure,
raw RL is often not enough on its own. Representation learning, shaping,
imitation, and task reformulation often matter just as much as the
choice of optimizer. Our project followed that pattern closely.

\subsection{Baselines and Experimental
Setup}\label{baselines-and-experimental-setup}

Our main baseline is the greedy spymaster described above. This is a
meaningful baseline because it already uses semantic information and is
much stronger than random clueing. In addition to that baseline, we
planned ablation experiments comparing SAC with HER but without reward
shaping or behavioral cloning, SAC with HER and reward shaping but
without behavioral cloning, and the full method with reward shaping,
behavioral cloning, and HER. These ablations are meant to test the
central hypothesis from the proposal: that semantic shaping and
imitation-based warm starts should improve learning over plain RL.

The current reproducible training setup uses a 5x5 board, 8 friendly
words, a goal size of 3, up to 8000 clue candidates, behavioral cloning
from 60 demonstration episodes, 3 BC epochs, and 8000 SAC training
timesteps. HER uses 4 sampled goals with the \texttt{future} selection
strategy. Our evaluation reports mean return, mean turns, win rate,
assassin rate, and friendly reveal rate. Those metrics were chosen
because win rate alone is not enough in this domain: an agent can be
conservative and avoid the assassin while revealing too few friendly
words, or it can reveal many friendly words but still fail through risky
clues.

To reproduce the current pipeline, one can create a Python 3.12
environment, install the project requirements, run the main SAC/HER
training configuration, and then run the ablation suite using the same
base settings. We also built an end-to-end smoke test to verify the full
pipeline before larger experiments.

\subsection{Experimental Settings
Table}\label{experimental-settings-table}

\begin{longtable}[]{@{}ll@{}}
\toprule\noalign{}
Setting & Value \\
\midrule\noalign{}
\endhead
\bottomrule\noalign{}
\endlastfoot
Board size & 5 x 5 \\
Friendly / Opponent / Neutral / Assassin & 8 / 8 / 8 / 1 \\
Max turns & 20 \\
Max clue count & 9 \\
Goal size & 3 \\
Embedding model & \texttt{all-MiniLM-L6-v2} \\
Clue vocabulary size cap & 8000 \\
BC enabled & Yes \\
Demonstration episodes & 60 \\
BC epochs & 3 \\
RL algorithm & SAC \\
HER enabled & Yes \\
Total timesteps & 8000 \\
Evaluation episodes in current base config & 3 \\
Seed in current base config & 43 \\
\end{longtable}

\subsection{Baseline and Ablation
Template}\label{baseline-and-ablation-template}

\begin{longtable}[]{@{}
  >{\raggedright\arraybackslash}p{(\columnwidth - 8\tabcolsep) * \real{0.2000}}
  >{\raggedright\arraybackslash}p{(\columnwidth - 8\tabcolsep) * \real{0.2000}}
  >{\raggedright\arraybackslash}p{(\columnwidth - 8\tabcolsep) * \real{0.2000}}
  >{\raggedright\arraybackslash}p{(\columnwidth - 8\tabcolsep) * \real{0.2000}}
  >{\raggedright\arraybackslash}p{(\columnwidth - 8\tabcolsep) * \real{0.2000}}@{}}
\toprule\noalign{}
\begin{minipage}[b]{\linewidth}\raggedright
Method
\end{minipage} & \begin{minipage}[b]{\linewidth}\raggedright
Reward shaping
\end{minipage} & \begin{minipage}[b]{\linewidth}\raggedright
Behavioral cloning
\end{minipage} & \begin{minipage}[b]{\linewidth}\raggedright
HER
\end{minipage} & \begin{minipage}[b]{\linewidth}\raggedright
Notes
\end{minipage} \\
\midrule\noalign{}
\endhead
\bottomrule\noalign{}
\endlastfoot
Greedy baseline & N/A & N/A & N/A & Non-learning semantic heuristic \\
SAC + HER & No & No & Yes & Plain RL-style baseline \\
SAC + HER + shaping & Yes & No & Yes & Tests the value of reward
shaping \\
Full method & Yes & Yes & Yes & Proposed final approach \\
Optional additional baseline & \texttt{TBD} & \texttt{TBD} &
\texttt{TBD} & Add if you include another comparison in the
presentation \\
\end{longtable}

\subsection{Results}\label{results}

Final experimental results are still being finalized, so the main
quantitative tables and claims are left as \texttt{TBD}. The intended
final comparison will report the greedy baseline, SAC/HER without
shaping or BC, SAC/HER with shaping only, and the full method. For each,
we plan to report mean return, win rate, assassin rate, friendly reveal
rate, mean turns, and the number of evaluation episodes.

Although those final results are not ready yet, we do have smoke-test
and debugging runs that confirm the pipeline works end to end. In those
small runs, the greedy baseline performs strongly while the short-run
BC-pretrained and SAC+HER models remain weak. That is exactly what we
would expect: those runs were designed to verify the full training loop,
not to produce a competitive final policy. We therefore treat them only
as sanity checks and not as evidence for the final claims of the
project.

\subsubsection{Main Results Table
Template}\label{main-results-table-template}

\begin{longtable}[]{@{}
  >{\raggedright\arraybackslash}p{(\columnwidth - 12\tabcolsep) * \real{0.1429}}
  >{\raggedright\arraybackslash}p{(\columnwidth - 12\tabcolsep) * \real{0.1429}}
  >{\raggedright\arraybackslash}p{(\columnwidth - 12\tabcolsep) * \real{0.1429}}
  >{\raggedright\arraybackslash}p{(\columnwidth - 12\tabcolsep) * \real{0.1429}}
  >{\raggedright\arraybackslash}p{(\columnwidth - 12\tabcolsep) * \real{0.1429}}
  >{\raggedright\arraybackslash}p{(\columnwidth - 12\tabcolsep) * \real{0.1429}}
  >{\raggedright\arraybackslash}p{(\columnwidth - 12\tabcolsep) * \real{0.1429}}@{}}
\toprule\noalign{}
\begin{minipage}[b]{\linewidth}\raggedright
Method
\end{minipage} & \begin{minipage}[b]{\linewidth}\raggedright
Mean return
\end{minipage} & \begin{minipage}[b]{\linewidth}\raggedright
Win rate
\end{minipage} & \begin{minipage}[b]{\linewidth}\raggedright
Assassin rate
\end{minipage} & \begin{minipage}[b]{\linewidth}\raggedright
Friendly reveal rate
\end{minipage} & \begin{minipage}[b]{\linewidth}\raggedright
Mean turns
\end{minipage} & \begin{minipage}[b]{\linewidth}\raggedright
Eval episodes
\end{minipage} \\
\midrule\noalign{}
\endhead
\bottomrule\noalign{}
\endlastfoot
Greedy baseline & \texttt{TBD} & \texttt{TBD} & \texttt{TBD} &
\texttt{TBD} & \texttt{TBD} & \texttt{TBD} \\
SAC + HER & \texttt{TBD} & \texttt{TBD} & \texttt{TBD} & \texttt{TBD} &
\texttt{TBD} & \texttt{TBD} \\
SAC + HER + shaping & \texttt{TBD} & \texttt{TBD} & \texttt{TBD} &
\texttt{TBD} & \texttt{TBD} & \texttt{TBD} \\
Full method & \texttt{TBD} & \texttt{TBD} & \texttt{TBD} & \texttt{TBD}
& \texttt{TBD} & \texttt{TBD} \\
\end{longtable}

\subsubsection{Optional Multi-Seed Summary
Table}\label{optional-multi-seed-summary-table}

\begin{longtable}[]{@{}
  >{\raggedright\arraybackslash}p{(\columnwidth - 10\tabcolsep) * \real{0.1667}}
  >{\raggedright\arraybackslash}p{(\columnwidth - 10\tabcolsep) * \real{0.1667}}
  >{\raggedright\arraybackslash}p{(\columnwidth - 10\tabcolsep) * \real{0.1667}}
  >{\raggedright\arraybackslash}p{(\columnwidth - 10\tabcolsep) * \real{0.1667}}
  >{\raggedright\arraybackslash}p{(\columnwidth - 10\tabcolsep) * \real{0.1667}}
  >{\raggedright\arraybackslash}p{(\columnwidth - 10\tabcolsep) * \real{0.1667}}@{}}
\toprule\noalign{}
\begin{minipage}[b]{\linewidth}\raggedright
Method
\end{minipage} & \begin{minipage}[b]{\linewidth}\raggedright
Seeds
\end{minipage} & \begin{minipage}[b]{\linewidth}\raggedright
Mean win rate
\end{minipage} & \begin{minipage}[b]{\linewidth}\raggedright
Std. dev. win rate
\end{minipage} & \begin{minipage}[b]{\linewidth}\raggedright
Mean return
\end{minipage} & \begin{minipage}[b]{\linewidth}\raggedright
Std. dev. return
\end{minipage} \\
\midrule\noalign{}
\endhead
\bottomrule\noalign{}
\endlastfoot
Greedy baseline & \texttt{TBD} & \texttt{TBD} & \texttt{TBD} &
\texttt{TBD} & \texttt{TBD} \\
SAC + HER & \texttt{TBD} & \texttt{TBD} & \texttt{TBD} & \texttt{TBD} &
\texttt{TBD} \\
SAC + HER + shaping & \texttt{TBD} & \texttt{TBD} & \texttt{TBD} &
\texttt{TBD} & \texttt{TBD} \\
Full method & \texttt{TBD} & \texttt{TBD} & \texttt{TBD} & \texttt{TBD}
& \texttt{TBD} \\
\end{longtable}

\subsubsection{Qualitative Analysis
Template}\label{qualitative-analysis-template}

\begin{longtable}[]{@{}
  >{\raggedright\arraybackslash}p{(\columnwidth - 6\tabcolsep) * \real{0.2500}}
  >{\raggedright\arraybackslash}p{(\columnwidth - 6\tabcolsep) * \real{0.2500}}
  >{\raggedright\arraybackslash}p{(\columnwidth - 6\tabcolsep) * \real{0.2500}}
  >{\raggedright\arraybackslash}p{(\columnwidth - 6\tabcolsep) * \real{0.2500}}@{}}
\toprule\noalign{}
\begin{minipage}[b]{\linewidth}\raggedright
Example board/clue situation
\end{minipage} & \begin{minipage}[b]{\linewidth}\raggedright
What the greedy baseline did
\end{minipage} & \begin{minipage}[b]{\linewidth}\raggedright
What the learned policy did
\end{minipage} & \begin{minipage}[b]{\linewidth}\raggedright
Interpretation
\end{minipage} \\
\midrule\noalign{}
\endhead
\bottomrule\noalign{}
\endlastfoot
\texttt{TBD} & \texttt{TBD} & \texttt{TBD} & \texttt{TBD} \\
\texttt{TBD} & \texttt{TBD} & \texttt{TBD} & \texttt{TBD} \\
\end{longtable}

\subsubsection{Result Discussion
Template}\label{result-discussion-template}

Once the final results are available, this section should answer four
concrete questions. First, did the full method outperform the plain RL
ablations? Second, how large was the gap between the learned policies
and the greedy baseline? Third, did reward shaping mainly improve sample
efficiency, final performance, or both? Fourth, did behavioral cloning
reduce instability or only provide a short-lived warm start? Writing the
discussion around these questions will make the final report much more
focused than simply restating the table entries.

\subsection{Deviations from the
Proposal}\label{deviations-from-the-proposal}

The final project stayed close to the proposal in spirit, but there were
several important deviations. First, we fixed the teammate guesser
instead of trying to learn both players jointly. We made this decision
because jointly training a spymaster and a guesser would have introduced
much more instability and would have made it harder to tell whether
gains came from better clueing or simply from a different teammate
policy. Fixing the guesser turned the project into a much cleaner
controlled experiment.

Second, we did not treat clue generation as unconstrained language
generation. Instead, the policy outputs a continuous embedding that is
decoded to the nearest legal clue in a fixed vocabulary. This was mainly
a practical decision. It keeps the action space manageable for SAC,
preserves valid game actions, and still allows the learned policy to
operate in a semantic space rather than a small hand-coded action set.

Third, the final system consolidated around sentence-transformer
embeddings rather than experimenting with a larger variety of symbolic
or lexical resources. This helped us keep the project coherent and
reproducible. Once the environment, baseline, and reward function were
all using the same semantic representation, the rest of the pipeline
became easier to build and debug.

Finally, the current training scale is still modest. Much of the project
effort went into getting a correct, inspectable, reproducible pipeline
rather than immediately scaling to long expensive runs. That is one
reason the final results are still pending.

Taken together, these deviations reflect a general theme of the project:
narrowing the scope was often necessary to make the core RL question
answerable. Instead of asking whether we could solve all of
\emph{Codenames} at once, we asked whether a spymaster can learn better
clueing behavior in a controlled semantic environment. That narrower
question was much more realistic for the class project timeline and
still scientifically meaningful.

\subsection{What We Learned}\label{what-we-learned}

One of the main things we learned is that this problem is much more
tractable once it is expressed in semantic space. A clue in Codenames is
not naturally a low-dimensional symbolic action. Thinking of clueing as
movement in an embedding space made several parts of the project line up
at once: the action space, the baseline, the reward shaping, and the
guesser model all became easier to define.

We also learned that sparse reward alone would likely be too weak for
this task. If the only signal is eventual success or failure, then most
early trajectories look equally bad. Reward shaping gave the agent
feedback about whether a clue was at least semantically promising, even
if the full turn or full game was not yet successful. That did not solve
the problem by itself, but it made the learning problem much better
structured.

Another lesson was that the greedy baseline was useful for more than
evaluation. It became an important tool for debugging and for generating
demonstrations. In a task like this, random exploration is unlikely to
stumble into good semantic clues often enough to be efficient.
Behavioral cloning gives the policy a much better starting point, and
even when the final learned policy does not beat the baseline, the
baseline still helps us build and understand the system.

More broadly, the experience reinforced a course-level lesson: combining
methods is often harder, but also more realistic, than using any one
method in isolation. Our project was not ``just SAC'' or ``just HER.''
The working version depended on a combination of representation design,
reward shaping, imitation, and off-policy RL. That layering was one of
the most important practical takeaways from the project.

There was also a more practical lesson about experimentation. In a
project like this, it is easy to confuse ``the method is weak'' with
``the pipeline is broken.'' Building small smoke tests and debugging
runs first turned out to be essential. That process did not make the
final learned policy stronger by itself, but it made our conclusions
more trustworthy because we could separate infrastructure issues from
actual learning difficulty.

\subsection{Development Path}\label{development-path}

The path to the current system was iterative rather than linear. Early
on, the main question was how to represent clues and board words in a
way that made the game learnable. Once we committed to embedding-based
semantics, it became natural to build a cosine-margin baseline, and that
baseline quickly became the center of the rest of the system. It
produced clues that were good enough to use as demonstrations, and it
also exposed when the clue vocabulary or similarity structure was not
behaving as expected.

A second important step was separating correctness from performance. We
built smoke-test configs and notebooks that verify environment reset and
step logic, clue decoding, demonstration generation, and short SAC/HER
runs. Those small tests were essential because they let us debug the
mechanics of the system without confusing infrastructure bugs with
genuine learning difficulty. In hindsight, that separation saved a lot
of time.

A third lesson from the development process was that simplifying the
game can make the scientific story stronger. Fixing the guesser,
constraining clues to a filtered vocabulary, and using coherent target
subsets all reduce realism somewhat, but they also make the problem
better defined and easier to evaluate. Those simplifications were not
accidental shortcuts; they were part of turning a very open-ended game
into a tractable RL problem.

If we were continuing the project beyond the class, the next stage would
be to scale experiments rather than redesign the method from scratch. At
this point the project already has a clear environment, a meaningful
baseline, and a structured RL pipeline. The remaining work is mostly
empirical: run longer training, evaluate across seeds, and determine
which combination of shaping, imitation, and HER actually matters most.

\subsection{Limitations and Future
Work}\label{limitations-and-future-work}

The current project still has several clear limitations. The guesser is
fixed rather than learned, clue generation is limited to a candidate
vocabulary, and the current training budgets are still relatively small.
The environment also approximates the real game's legal clue rules
rather than modeling every edge case. Because of those choices, the
project should be understood as a structured RL study of the spymaster
role, not as a complete superhuman Codenames system.

The most obvious next steps are to scale up training, evaluate across
multiple random seeds, test more embedding models, and complete the
ablation study with stronger final runs. A natural extension would be to
compare the current continuous embedding-action approach against a
discrete-action formulation over clue vocabularies, or to replace the
fixed guesser with stronger or more human-like teammate policies.

\subsection{Reproducibility Checklist}\label{reproducibility-checklist}

\begin{longtable}[]{@{}
  >{\raggedright\arraybackslash}p{(\columnwidth - 4\tabcolsep) * \real{0.3333}}
  >{\raggedright\arraybackslash}p{(\columnwidth - 4\tabcolsep) * \real{0.3333}}
  >{\raggedright\arraybackslash}p{(\columnwidth - 4\tabcolsep) * \real{0.3333}}@{}}
\toprule\noalign{}
\begin{minipage}[b]{\linewidth}\raggedright
Item
\end{minipage} & \begin{minipage}[b]{\linewidth}\raggedright
Included in project?
\end{minipage} & \begin{minipage}[b]{\linewidth}\raggedright
Notes
\end{minipage} \\
\midrule\noalign{}
\endhead
\bottomrule\noalign{}
\endlastfoot
Fixed board vocabulary & Yes & Same board word source across runs \\
Fixed clue vocabulary construction & Yes & Filtered and capped clue
set \\
Configurable hyperparameters & Yes & Main settings stored in config
files \\
Deterministic seeds & Yes & Seeded runs supported \\
Baseline implementation & Yes & Greedy semantic baseline available \\
Ablation setup & Yes & Core variants already defined \\
Smoke-test pipeline & Yes & End-to-end debugging workflow included \\
Final large-scale results & \texttt{TBD} & Still running / not
finalized \\
\end{longtable}

\subsection{Conclusion}\label{conclusion}

This project implemented a full RL pipeline for the Codenames spymaster
task, grounded in the ideas from our proposal. The final system includes
a custom goal-conditioned environment, semantic embeddings for board
words and clues, a greedy semantic baseline, behavioral cloning from
demonstrations, and SAC training with optional HER and ablations. Even
before the final result runs are complete, the project already provides
a reproducible framework for studying semantic decision-making in a
word-based game.

The biggest conceptual contribution of the project is not any single
algorithmic component. It is the way the task was reformulated so that
reinforcement learning becomes plausible at all. By treating clueing as
semantic control in embedding space, and by combining reward shaping
with imitation and HER, we turned a difficult language-like problem into
something that can be explored with the RL tools from class.

\subsection{Bibliography}\label{bibliography}

Tu, I., and teammates. \emph{Codenames Proposal}. Project proposal PDF
included in the repository as \texttt{Codenames\ Proposal.pdf}.

Haarnoja, T., Zhou, A., Abbeel, P., and Levine, S. ``Soft Actor-Critic:
Off-Policy Maximum Entropy Deep Reinforcement Learning with a Stochastic
Actor.'' \emph{Proceedings of ICML}, 2018.

Andrychowicz, M., Wolski, F., Ray, A., et al.~``Hindsight Experience
Replay.'' \emph{Proceedings of NeurIPS}, 2017.

Reimers, N., and Gurevych, I. ``Sentence-BERT: Sentence Embeddings using
Siamese BERT-Networks.'' \emph{Proceedings of EMNLP-IJCNLP}, 2019.

Stable-Baselines3 documentation and source code.

Gymnasium documentation.

PyTorch documentation.

Sentence-Transformers documentation.

Chvatil, Vlaada. \emph{Codenames}. Czech Games Edition, 2015.

\subsection{Documentation of Sources Used for Code, Including
LLMs}\label{documentation-of-sources-used-for-code-including-llms}

The implementation relies on external libraries rather than copied
standalone code. In particular, Stable-Baselines3 was used for SAC and
HER support, Gymnasium for the environment interface, PyTorch for
supervised actor pretraining, and Sentence-Transformers for semantic
embeddings. The project also used general library documentation for
those packages during development.

LLM assistance was used during the project as a coding and writing aid.
This included help with implementation scaffolding, debugging
suggestions, code explanation, refactoring ideas, and drafting the
report. All final design choices, integration decisions, and submitted
conclusions were reviewed and edited by the project team.

\end{document}
